% !TEX root = ../../../main.tex
% Week 4, IMG_9234--IMG_9235 and IMG_9238.
\section{Number series}
\label{sec:analysis-i:number-series}

An infinite sum is defined through finite sums. This lets us use the
sequence tools from the preceding chapter.

\begin{definition}
  Given real numbers $a_1,a_2,\ldots$, their \emph{partial sums} are
  $S_n=\sum_{k=1}^n a_k$. The series $\sum_{k=1}^{\infty}a_k$
  \emph{converges} to $S\in\R$ if $S_n\to S$. If the partial sums
  have no finite limit, the series \emph{diverges}.
\end{definition}

\begin{theorem}[Cauchy criterion for series]
  \label{thm:analysis-i:series-cauchy}
  The series $\sum_{k=1}^{\infty}a_k$ converges if and only if
  \[
    \forall\varepsilon>0\ \exists N\in\N\quad
    m>n\ge N\ \Longrightarrow\
    \left|\sum_{k=n+1}^{m}a_k\right|<\varepsilon.
  \]
\end{theorem}

\begin{proof}
  For $m>n$, split $S_m$ after its first $n$ terms and subtract $S_n$.
  Using the definition of the partial sums gives
  \begin{align*}
    S_m-S_n
      &=\sum_{k=1}^{m}a_k-\sum_{k=1}^{n}a_k\\
      &=\left(\sum_{k=1}^{n}a_k+\sum_{k=n+1}^{m}a_k\right)
        -\sum_{k=1}^{n}a_k\\
      &=\sum_{k=n+1}^{m}a_k.
  \end{align*}
  The shared terms $a_1,\ldots,a_n$ cancel. The first remaining term
  is $a_{n+1}$ and the last is $a_m$, so the remaining sum has $m-n$
  terms. In particular, when $m=n+1$, it contains just $a_{n+1}$.
  Taking absolute values therefore gives
  \[
    |S_m-S_n|=\left|\sum_{k=n+1}^{m}a_k\right|.
  \]
  Thus the expression in the criterion measures the distance between
  two partial sums, using precisely the terms added after the $n$th one.

  The restriction $m>n$ simply orders the two partial sums. Exchanging the indices
  gives the same absolute difference, and equal indices give zero.
  Thus checking $m>n\ge N$ is equivalent to checking every pair
  $m,n\ge N$. The criterion therefore says exactly that $(S_n)$
  is Cauchy. By \autoref{thm:analysis-i:cauchy-completeness}, this is
  equivalent to convergence of $(S_n)$ in $\R$, which is the definition
  of convergence of the series.
\end{proof}

\begin{proposition}
  \label{thm:analysis-i:series-term-test}
  If $\sum_{n=1}^{\infty}a_n$ converges, then $a_n\to0$.
\end{proposition}

\begin{proof}
  If $S_n\to S$, then $S_{n-1}\to S$ as well. Thus
  $a_n=S_n-S_{n-1}\to S-S=0$ by the difference rule for limits.
\end{proof}

This condition is necessary but is not sufficient. The harmonic series
provides the basic counterexample.

\subsection{The harmonic series}

Let $H_n=\sum_{k=1}^n1/k$. Although $1/n\to0$, for every $N\ge1$
we have
\[
  H_{2N}-H_N=\sum_{k=N+1}^{2N}\frac1k
  \ge N\frac{1}{2N}=\frac12.
\]
There are $N$ terms in this block, each at least $1/(2N)$.
Hence $(H_n)$ fails the Cauchy condition with $\varepsilon=1/2$,
so the harmonic series diverges. Its partial sums are increasing;
if they were bounded above, the monotone bounded sequence theorem
would force convergence. They are therefore unbounded and tend to
$+\infty$. \autoref{fig:harmonic-blocks} shows why each doubling
contributes at least another $1/2$.
% IMG_9235 writes a strict bound; >= includes N=1.

\begin{figure*}[htbp]
  \centering
  % !TEX root = ../../main.tex
% Shared physical dimensions for interval and number-line diagrams.
\tikzset{
  interval diagram/.style={
    x=1cm,y=6mm,font=\footnotesize,
    color=black,line width=0.5pt,line cap=round,line join=round,
    >={Stealth[length=4pt,width=3pt]}
  },
  interval endpoint/.style={
    circle,draw=black,line width=0.5pt,
    minimum size=4pt,inner sep=0pt,outer sep=0pt
  },
  interval open/.style={interval endpoint,fill=white},
  interval closed/.style={interval endpoint,fill=black},
  interval label/.style={inner sep=1pt}
}

\begin{tikzpicture}[interval diagram,x=6mm,y=45mm]
  \foreach \start/\count/\height/\heightlabel in
    {0/1/0.5/\frac12,3/2/0.25/\frac14,7/4/0.125/\frac18,13/8/0.0625/\frac1{16}}{
    \fill[black!8] (\start,0) rectangle ({\start+\count},\height);
    \draw (\start,0) rectangle ({\start+\count},\height);
    \foreach \i in {1,...,\count}{
      \draw ({\start+\i},0) -- ({\start+\i},\height);
    }
    \node[above=6pt] at ({\start+\count/2},\height) {$N=\count$};
    \node[below=8pt] at ({\start+\count/2},0) {$\heightlabel$};
  }
\end{tikzpicture}

  \caption[Harmonic blocks with a constant lower bound.]{Lower bounds for four harmonic blocks. A block of $N$ terms
    has $N$ strips of height $1/(2N)$, hence area $1/2$. The actual
    terms are at least these heights, so the block sums do not shrink.}
  \label{fig:harmonic-blocks}
\end{figure*}

\subsection{The geometric series}

For $x\ne1$, multiplying a finite sum by $1-x$ gives
\[
  (1-x)\sum_{k=0}^N x^k=1-x^{N+1},\qquad
  \sum_{k=0}^N x^k=\frac{1-x^{N+1}}{1-x}.
\]
If $|x|<1$, then $|x|^{N+1}\to0$ by the geometric-sequence
estimate in \autoref{sec:analysis-i:convergent-sequences}, with
$x=0$ immediate. Consequently
\[
  \sum_{k=0}^{\infty}x^k=\frac{1}{1-x}\qquad(|x|<1).
\]
If $|x|\ge1$, the terms $x^k$ do not tend to zero, so the series
diverges by \autoref{thm:analysis-i:series-term-test}.
% IMG_9238's repeated claim a^n -> 1 for 0<a<1 is a slip; it is 0.

The Cauchy criterion can also establish convergence when the terms
alternate in sign.

% !TEX root = ../../hw04.tex
% Source: HW4 Intro to Math Analysis I.pdf, page 1, question 3.
\begin{exercise}
  \label{ex:hw04:alternating-series}
  \solutionlink{hw04:alternating-series}
  An alternating series is one of the following form
  \[
    \sum_{n=0}^{\infty}(-1)^n a_n,
    \qquad
    \begin{gathered}
      a_n\ge0,\quad a_n\ge a_{n+1},\quad n=0,1,2,\ldots;\\
      \lim_{n\to\infty}a_n=0.
    \end{gathered}
  \]
  Prove that an alternating series converges.
  (Hint: Use the Cauchy convergence criterion.)
\end{exercise}

